
\documentclass[11pt]{article}
\usepackage[a4paper,margin=1.05in]{geometry}
\usepackage[T1]{fontenc}
\usepackage[utf8]{inputenc}
\usepackage{lmodern}
\usepackage{microtype}
\usepackage{amsmath,amssymb}
\usepackage{setspace}
\usepackage{hyperref}
\hypersetup{colorlinks=false,pdfborder={0 0 0}}
\setlength{\parindent}{1.15em}
\setlength{\parskip}{0.18em}
\sloppy
\begin{document}
\begin{center}
{\Large\bfseries Memorial Address on Gustav Peter Lejeune Dirichlet\par}
\vspace{0.35em}
{\large E. E. Kummer\par}
\vspace{0.35em}
{\small Read at the public session of the Royal Academy of Sciences on 5 July 1860.}\par
\vspace{0.4em}
{\small Source: G. Lejeune Dirichlet, \emph{Werke}, Band II, pp. 309--344; appendix item XXVII.}\par
\vspace{1.2em}
\end{center}
It is not yet ten years since the three men to whom our German fatherland owes a new flowering of the mathematical sciences - Gauss, Jacobi, and Dirichlet - were still alive and still active, working to renew and confirm, with new splendour, the old fame for deep knowledge of the most abstract mathematical truths, as well as of those mathematical truths concretely realized in nature, which Kepler and Leibniz above all had won for the German nation. At that time our Academy had the good fortune to possess two of these outstanding men as active members, Jacobi and Dirichlet. Personally bound by friendship, they stimulated and advanced one another by the free communication of their deep mathematical thoughts, and exercised the most lasting influence on the general development of the mathematical sciences. Jacobi's early death was the first irreparable loss suffered by the science that had unfolded to such bloom in our country. The importance of the creations of that rarely gifted investigator, and the eminent position which he will occupy for all time in the history of mathematics, were described by Dirichlet in the memorial address delivered here eight years ago today, so profoundly and so truly that he thereby erected the fairest and worthiest monument to the memory of the departed. When, four years after Jacobi, the aged Gauss passed from life in the undisputed fame of being the first mathematician of his time, this great general loss had for our Academy the further deplorable consequence that Dirichlet, as the only worthy successor to that great man, was called to Göttingen and left the number of its resident members.

The Academy, unable either to avert or to replace this loss, preserved by electing him one of its ordinary non-resident members the right still to regard Dirichlet as its own. For his special colleagues he remained the living centre of their investigations and labours until death set an end to his life and work. Our Academy, to which Dirichlet belonged for twenty-seven years, and in whose writings his imperishable masterpieces are deposited, has the right to regard the scientific fame of this great mathematician as its own; and above all others it has therefore the duty to preserve his memory and to pay him the last academic honour in a public memorial address. The veneration which I myself have always felt for the departed, the friendship which he bestowed on me and preserved for more than twenty years, and the close relation in which my own studies stand to his scientific work, have moved me to take the floor here before this highly distinguished assembly, in order to describe the great scientific significance of his masterpieces and at the same time, in a few strokes, to sketch the image of his life and of his character, which was noble and pure like his writings. I know that I shall be able to solve the task set before me only very imperfectly: not only for the substantive reason that the true significance of the intellectual creations of great men can often be properly recognized and appreciated only in the later historical course of science, where they not seldom become starting-points of richly unfolding theories, but also because of my own weakness, for which I must ask your kind indulgence.

Gustav Peter Lejeune Dirichlet was born on 13 February 1805 in Düren. His father, who held there the office of post director, a gentle, obliging, and lovable man, and his mother, still living now at a very advanced age, a spirited and finely cultivated woman, gave the boy, endowed by nature with more than ordinary gifts, a very careful upbringing. He received his first instruction in an elementary school; when this was no longer found sufficient for him, he attended a private school, and in addition received special instruction in Latin so that he might later attend a gymnasium. His great preference for mathematics appeared very early. When he was not yet twelve years old, he used his pocket money to buy mathematical books, with which he occupied himself very diligently, especially in the evenings. When people told him that he could not possibly understand them, he answered: I read them until I understand them. His parents wished him to become a merchant; but when he showed decided aversion to that path, they yielded and in 1817 sent him to the gymnasium at Bonn.

Through the friendly communication of Professor Elvenich of Breslau, who at that time, as a student in Bonn, lived in the same house as the young Dirichlet and had been urgently recommended by the anxious parents to supervise and guide him, I am able to give the following faithful and lively account of the boy, then about thirteen years old. In his behaviour he distinguished himself very favourably by propriety and good manners; and the unselfconsciousness and openness of his whole nature made all who had to do with him sincerely well disposed toward him. His diligence was orderly, yet directed chiefly to mathematics and history. He studied even when he had no schoolwork to do, for even then his active mind was always occupied with worthy objects of reflection. Great historical events, notably the French Revolution, and public affairs interested him in a high degree; and he judged these and other matters with an independence unusual for his youth, from the standpoint of a liberal cast of thought, perhaps the fruit of his upbringing at home. Everything coarse and ignoble was always repugnant to him. Games and other youthful amusements, however, had almost no attraction for him, while he loved sociable conversation and took part gladly and animatedly, especially in discussions of politics and historical subjects. In general his mind, whose prominent quality was sharp discernment, moved in a much higher sphere than is usual at that age.

Dirichlet remained at the Bonn gymnasium only two years and then exchanged it for the Jesuit Gymnasium in Cologne, which his parents preferred for reasons unknown to me. Here his teacher in mathematics was Georg Simon Ohm, later famous for the discovery of the law of electrical resistance that bears his name. Under Ohm's instruction, and by his own diligent study of mathematical works, Dirichlet made very considerable progress in that science and acquired an unusual breadth of knowledge. Yet he by no means neglected the other disciplines, and he passed through the course of the gymnasium very rapidly, so that as early as 1821, when he was only sixteen, he received the leaving certificate for the university and returned home to discuss with his parents the choice of his future vocation. It was natural that they should meet his decision to study mathematics with earnest admonition to secure his livelihood in the world by a more practical study, for which they proposed jurisprudence. He replied to them modestly but firmly that, if they required it, he would obey them, but that he could not abandon his favourite study and would at least devote the nights to it. The parents, as sensible as they were tender, then yielded to their son's decided wish.

Mathematical study at the Prussian and other German universities was then in a sorry state. The lectures, which rose only a little above the domain of elementary mathematics, were in no way fitted to satisfy the urge toward deeper knowledge which animated the young Dirichlet; and apart from the one great name of Gauss there was in Germany no other that could have exerted a special attraction on him. In France, by contrast, and especially in Paris, mathematics still stood in its full bloom. There a circle of men whose great names will shine for all time in the history of the mathematical sciences worked, in research and teaching, at their living development and dissemination. There still lived the great Laplace, to whom his Celestial Mechanics secured undisputed first rank, and who was still working on the completion of that work and on a supplement to his theory of probability. Legendre, restlessly active into old age, was perfecting his theory of elliptic functions by the discovery of a new transformation of them and preparing the third edition of his work on number theory. Fourier, who had recently completed his mathematical theory of heat, gathered around him a select circle of the most talented young mathematicians for scientific and cheerful conversations. Poisson enriched mechanics and mathematical physics by a series of the most valuable memoirs. Cauchy was then laying the foundation for an essential improvement and transformation of the whole of analysis, by stricter methods and by the introduction of imaginary variables. These men, and besides them a considerable number of other scientific notabilities, some of whom are still living, worked together to make Paris the most brilliant seat of the mathematical sciences.

Correctly appreciating these circumstances, Dirichlet recognized that Paris was the place where he could expect the greatest gain for his mathematical studies. Since his parents, who from earlier times still had connections with some friendly families in Paris, readily consented, in May 1822 he entered this high school of the mathematical sciences, with the joyful consciousness that he could now devote himself entirely to his favourite study. He attended lectures at the Collège de France and at the Faculté des Sciences, where he had Lacroix, Biot, Hachette, and Francoeur as teachers. An attempt to obtain permission to attend lectures at the École Polytechnique as a guest failed because the Prussian chargé d'affaires in Paris, without special authorization from Minister von Altenstein, would not undertake to obtain permission from the French ministry.

In addition to hearing lectures and thinking through the material presented in them, Dirichlet devoted his time to the strenuous study of the foremost mathematical writings, and among these above all to Gauss's work on higher arithmetic, the Disquisitiones Arithmeticae. This work exercised a much more important influence on his whole mathematical education and direction than his other Parisian studies. He not only studied it through once or several times; throughout his life he never ceased, by repeated reading, to make again present to himself the abundance of deep mathematical thoughts contained in it. For that reason it was never set up on his bookshelf, but had its permanent place on the table at which he worked. One can infer how much effort it must have cost him to work his way into this extraordinary book from the fact that more than twenty years after it appeared none of the living mathematicians of the time had completely studied and understood it; and that even Legendre, who had devoted a large part of his life to higher arithmetic, had to confess in the second edition of his number theory that he would gladly have enriched his work with Gauss's results, but that the methods of that author were so peculiar that he could not have reproduced them without the greatest detours or without taking on the role of a mere translator. Dirichlet was the first not only to understand this work completely, but also to open it up for others, making fluid and transparent the rigid methods behind which deep thoughts lay concealed, and in many main points replacing them by simpler, more genetic ones, without yielding the least in the perfect rigour of the proofs. He was also the first, going beyond it, to reveal a rich treasure of still deeper secrets of number theory.

Dirichlet's outward life in the first year of his stay in Paris was extremely simple and withdrawn. His studies, interrupted only once by an attack of smallpox, occupied him completely, and his social intercourse was limited to a few houses to which he had been recommended and to a few young Germans who were there for their studies. In the summer of 1823, however, a change occurred that was of the greatest importance for his general cultivation. General Foy, a man of spirit and many-sided education, distinguished no less by the eminent position he held as head of the opposition in the Chamber of Deputies and as one of its most celebrated orators than by his brilliant military career, and whose house was one of the most respected and sought after in Paris, was then looking for a young man as teacher for his children, chiefly to instruct them in German language and literature. Through the mediation of a friend of the Dirichlet family, M. Larchet de Chamont, our Dirichlet was recommended to him. At their first personal meeting, the young man's open and modest manner made so favourable an impression on the general that he immediately offered him the position of teacher to his children, with a respectable salary and with such slight duties that he still had enough free time to continue the studies he had begun.

Dirichlet accepted with joy, not only because he was thereby placed in a position no longer to cost his parents anything, but especially because he expected much benefit for his outward worldly education, in which by his own judgment he was still very backward, from residence in the house of so broadly cultivated and distinguished a man. In this new position he felt extraordinarily content and happy, since Monsieur and Madame Foy showed him everywhere the greatest friendliness and courtesy and regarded him as a member of their own family. The instruction of the children, the eldest of whom was an eleven-year-old girl, cost him little effort; and the general's wife, who under his guidance zealously and with the best success took up again the German she had practised in childhood but had since completely forgotten, repaid his effort in an equally pleasant and useful way by helping him, through exercises in reading French, to lose the foreign accent in his pronunciation. The general, however, exercised the greatest influence on him, by the living example of an energetic, noble, and finely cultivated man; and this influence extended not merely to Dirichlet's outward cultivation, his habits and inclinations, but also to his way of thinking and acting and to his general views of life. It was also of great importance for his whole life that the general's house, a meeting-point of the chief notabilities in art and science of the French capital, and a place where the great political questions were discussed by the most respected members of the Chamber, questions whose nearest and provisional solution came in 1830, first gave him the opportunity to see life on a large scale and to take part in it.

All these new impressions that he absorbed, and the occupations and diversions connected with his position, did not turn Dirichlet away from his mathematical studies. On the contrary, precisely in this period he worked with intense diligence on his first published paper, Mémoire sur l'impossibilité de quelques équations indéterminées du cinquième degré. The subject of this memoir is most closely related to Fermat's assertion that the sum of two powers with the same exponent can never be equal to a power with the same exponent when these powers are higher than the second. This theorem, of which Fermat stated that he could prove it, but which more than 150 years after Fermat, at the time when Dirichlet occupied himself with it, had despite the strenuous efforts of Euler and Legendre been proved no farther than for the third and fourth powers, cannot indeed claim special value as an isolated item taken out of its scientific context. But it had acquired an unusually high importance because, as an apparently near yet distant goal set up in the then still quite unknown region of forms of higher degree, it had a most decisive influence on the direction taken by number theory in its historical development. Pursuing in his work the direction indicated by Fermat's theorem, Dirichlet considered the sum of two fifth powers, about which nothing had yet been established, and set himself the problem somewhat more generally: to investigate in what cases such a sum could not be equal to a given multiple of a fifth power. He levelled and secured the path of the investigation by some new theorems on the most general solution of the problem of making a quadratic form equal to a power number, and arrived at the proof of the impossibility of a whole class of equations of the fifth degree. Fermat's equation for fifth powers, one of which would necessarily have to be divisible by five, is included in this class only in the case where that one is also even; the other case, where it is odd, was shortly afterwards likewise shown to be impossible by Legendre, who followed a little further the path opened by Dirichlet. Thus Dirichlet shares with Legendre the honour of having advanced by a whole stage the proof of this historically remarkable theorem.

Not only the new results gained in one of the most difficult parts of number theory, but also the concision and sharpness of the proofs and the exceptional clarity of the exposition secured this first work of Dirichlet a brilliant success. The Paris Academy, to which he presented it, allowed him to read it at the session of 11 June 1825; and already at the next session, on the 18th of the same month, Messrs. Lacroix and Legendre, as commissioners, gave such a favourable report that the Academy resolved to include it in the collection of memoirs of foreign scholars. Dirichlet's reputation as an outstanding mathematician was thereby established; and as a young man from whom a great future could be expected he was thereafter not only admitted into the highest scientific circles of Paris but sought out by them. He thereby also entered into closer relations with several of the most respected members of the Paris Academy. Among them two must especially be mentioned: Fourier, who exerted an influence on the direction of his scientific researches, and Alexander von Humboldt, who exerted a significant influence on the later shaping of his outward life.

Fourier, who from the time of his youth, when he had taken an active part in the founding of the École Normale and the École Polytechnique, still preserved undiminished his enthusiasm for living scientific communication, and for whom it was an inner need to communicate orally what he had beautifully and greatly discovered, found in Dirichlet a young man to whom he could entirely open his mathematical heart and by whom he was not merely admired but also completely understood. He therefore drew him into the circle of distinguished young mathematicians whom he was accustomed to gather around himself, and with whom he then discussed his theory of heat, the new analytical methods invented for it, and various more general scientific subjects and questions in his own lively and attractive manner. In this circle, to which among others Sturm belonged, who shortly afterwards became generally famous through his theorem on the roots of algebraic equations, Dirichlet received manifold stimulation. In particular, Fourier awakened his interest in mathematical physics, in which he later worked with significant success; likewise Fourier series and integrals, which first received their true scientific foundation through Dirichlet's rigorous methods, occupy a not insignificant place in his later works.

Alexander von Humboldt, who was then living in Paris, had earlier heard General Foy praise Dirichlet as an outstanding mathematician, yet without attaching special weight to this praise, since it did not come from the mouth of a specialist. Only when Dirichlet, in consequence of the favourable reception that the Academy had granted to his memoir, paid his visit to Humboldt, did he become more closely acquainted with him. Humboldt received him with the most distinguished kindness and courtesy, and together with respect for his talent and scientific competence bestowed on him the liveliest personal sympathy and affection, which from then on he continually preserved and actively showed. Already during this first visit Dirichlet, in the course of conversation, let it be known that he intended later to return to his fatherland; Humboldt, who seized upon this thought with joy, strengthened him in his resolve by assuring him that there, with the extremely small number of distinguished German mathematicians, he could not fail, whenever he desired it, to obtain a very good position. Under the circumstances then prevailing, when negotiations for Gauss's call to Berlin, continued for several years, had just had to be abandoned because a few hundred thalers were lacking, this assurance was not so easy to fulfil; and soon afterwards Humboldt's tireless activity and all his influence were required to make it even approximately true.

The death in November 1825 of his highly honoured patron General Foy, and the influence of Alexander von Humboldt, who soon afterwards left Paris and moved to Berlin, brought Dirichlet to a mature decision to return to his fatherland. He addressed to Minister von Altenstein a request for a position suitable for him, which Humboldt undertook to support and to make effective by his influence; and in the autumn of 1826 he returned to his parents in Düren to await the result there.

While he was working there on a new memoir, to which I shall shortly return, Humboldt pursued the matter of his appointment with the liveliest zeal. He intervened personally with Minister von Altenstein to obtain for Dirichlet an extraordinary professorship at a Prussian university, with a salary of six to eight hundred thalers; he drew the most respected members of our Academy into his interest, so that their recommendations might support his own; and he frequently sent Dirichlet reports and good advice about what he should do on his side. But through all these efforts, even supported by Gauss through a letter addressed to our colleague Encke and handed by him to the Royal Ministry, no more could be attained than the assurance of 400 thalers a year as a fixed remuneration, so that he might habilitate at Breslau as a Privatdozent. Since the fixed remuneration secured him for the time being a modest livelihood, and since he could rely on Humboldt to continue his efforts to obtain for him a more suitable position, he accepted without hesitation. Meanwhile the philosophical faculty of the University of Bonn had created him Doctor of Philosophy honoris causa, which substantially eased his habilitation at a university.

On his journey to Breslau he chose the route through Göttingen in order to make the personal acquaintance of Gauss, and visited him on 18 March 1827. I have not been able to ascertain more detailed information about this meeting. A letter addressed to his mother from that time says only that Gauss received him very kindly and that the personal impression made by this great man was much more favourable than he had expected.

In Breslau, according to the statutes of the philosophical faculty there, he was to obtain the venia docendi by holding a trial lecture and colloquium before the faculty, writing a dissertation, and publicly defending it in Latin. For these performances, insofar as they concerned his science, he was more than equal to the task. He also knew very well how to write clearly and correctly in Latin on a scientific subject. But he had never devoted his time, consecrated to higher scientific ends, to acquiring the outward facility of speaking Latin; the empty formality of the Latin disputation was therefore highly disturbing and unpleasant to him. To his great satisfaction, but to the great regret of some gentlemen of the faculty there, after he had given his trial lecture on the irrationality of the number pi, the Royal Ministry entirely dispensed him from the public disputation. And so that he might immediately begin his lectures, without which the subject of mathematics there would have been almost wholly unrepresented, he received at the same time permission to submit his Latin habilitation dissertation afterwards.

The success of his teaching activity during the three semesters in which he lectured in Breslau was not significant. The students there, unwilling to go beyond the narrow circle of mathematical ideas and thoughts that had previously been transmitted to them in lectures, could not easily accustom themselves to his manner of teaching, which was unfamiliar to them; and his modest, even somewhat shy, appearance was not suited to impose on them. In Breslau generally, Dirichlet was gladly seen and sought in all social circles as a finely cultivated and lovable young man, but as a mathematician he was held in low esteem in comparison with his predecessor, who had written a textbook of analytic geometry. Since he himself never spoke of himself or of his own scientific merits, and had no literary following to do this for him, he did not attain there that local or provincial fame which in more restricted circles is more effective than the general recognition of the first men of science.

During his Breslau stay Dirichlet wrote two papers, both prompted by Gauss's memoir on biquadratic residues. In April 1825 the Göttinger gelehrte Anzeigen had brought a brief announcement of a series of papers on biquadratic residues and their reciprocity laws which Gauss intended to publish. The first had already been presented to the Göttingen Society of Sciences, but appeared only three years later. This announcement, giving some of Gauss's results whose proofs were to rest on a wholly new principle of number theory, aroused in the highest degree the curiosity of both Jacobi and Dirichlet. Each sought, by a quite different path, to penetrate Gauss's secret, and each succeeded in finding in this domain of higher power residues an abundance of new theorems, although the new principle, consisting in the introduction of complex integers, still remained concealed from them. For the already published Gaussian theorems, which contained the complete solution of the problem of specifying all primes for which the number two is a biquadratic residue or nonresidue, Dirichlet found astonishingly simple proofs by the known methods of number theory. In the same way he also solved the more general question for all arbitrarily given primes, so that only one step remained to the complete reciprocity law for biquadratic residues, a step that became possible only through the new Gaussian principle. In the other work then published, which Dirichlet wrote in Latin and submitted to the philosophical faculty as his habilitation dissertation, he gave an example, then quite new, of forms of arbitrary high degree whose divisors have certain linear forms. The results of this work may now be regarded as special cases of the general theorems on the divisors of the norm forms of the complex numbers formed from roots of unity.

To give a measure of the value then attributed especially to the first of these two writings, I cite the judgments of Bessel and Fourier upon it. Bessel wrote of it in a letter to Humboldt: who would have thought that genius would succeed in reducing something that seemed so difficult to such simple considerations? The name Lagrange could stand over the paper and no one would notice the incorrectness. Fourier, however, placed Dirichlet's achievement even higher than the great discoveries of Jacobi and Abel in the theory of elliptic functions, of which, to be sure, he had until then had knowledge only through Legendre's praises, which he considered exaggerated. In a letter addressed to Dirichlet, and in oral remarks made to the above-mentioned friend of the Dirichlet family, M. Larchet de Chamont, from which this judgment is taken, he also expressed the lively wish that Dirichlet might return to Paris, because he was called to occupy very soon one of the first places in the Academy there.

Meanwhile Alexander von Humboldt had obtained Dirichlet's appointment as extraordinary professor at the University of Breslau and was now working to win him for the university here and for the Academy, but first of all to draw him to Berlin. Since a mathematical teaching position about to become free at the General War School offered the appropriate opportunity, Humboldt seized it and recommended Dirichlet most urgently for it to General von Radowitz and to the Minister of War. They could not, however, decide at once to entrust the position to him definitively, probably because at the time, being only twenty-three years old, he may have seemed to them too young for it. It was therefore arranged with Minister von Altenstein that Dirichlet should at first receive a year's leave in order to take over the teaching at the War School provisionally. In the autumn of 1828 he came to Berlin to assume this new position.

The mathematical lectures which he had to give there before officers of about his own age gave him much pleasure. Association with cultivated military men, to which he had been accustomed from his time in General Foy's house, suited him very well; and since during that period he had also made thorough studies in modern military history, this common interest connected him with his listeners outside mathematics as well. Only later, after he had formed at the university here a large circle of listeners who followed him with lively scientific interest into the highest regions of mathematics, where he preferred to move, did the wish arise in him to be released from teaching at the War School; this wish then became one of the chief motives for his move to Göttingen.

Soon after his arrival in Berlin Dirichlet also took the necessary steps to be allowed to lecture at the university here. As professor at another university he did not have that right; nothing remained to him but to habilitate once more as Privatdozent, and in this sense he addressed his petition to the philosophical faculty. The faculty, however, dispensed him from the habilitation requirements in view of his otherwise proven scientific competence; and so he initially lectured here under the legal title of a Privatdozent. His definitive transfer as extraordinary professor to the university here took place only in 1831; a few months later he was elected by our Academy to ordinary membership. In the same year he married Rebecca Mendelssohn-Bartholdy, a granddaughter of Moses Mendelssohn. Remarkably, Alexander von Humboldt himself had an involuntary share even in this, insofar as it was he who first introduced Dirichlet into the famous house of his parents-in-law, distinguished by spirit and artistic sense.

For a longer time the further events of his life now recede before the significance of the scientific works that Dirichlet produced during the twenty-seven years of his life here. In describing them, which is now my task, I shall try to present in larger and more general outlines the advances of science contained in them, considering them not individually by the time of their origin but grouped according to their contents and to the underlying ideas.

Dirichlet's purely analytical works on infinite series and definite integrals, and on the functions appearing in these forms, originally arose from his study of mathematical physics and especially of Fourier's theory of heat. In applying these general forms to physical problems, however, he could not content himself with using them as ready-made tools; and since closer examination soon showed him that even in the most important points they still lacked strict scientific foundation, he directed his work first to securing these foundations. The series progressing according to the sines and cosines of multiples of an arc, which Fourier had used with the most outstanding success to represent arbitrary functions occurring in the theory of heat, had up to then, in all cases where the function to be expanded does not become infinite, shown the remarkable property of always being convergent. Yet even Cauchy's efforts had not succeeded in proving this generally and rigorously. The path taken here by that mathematician, famous not less for the rigour than for the originality of his methods, namely to examine only the relative magnitudes of the individual terms of these series and to base his conclusions on that, did not lead to true knowledge of this hidden property, but only quite close to it; for the convergence of these series in certain cases also depends on the particular way in which their positive and negative terms cancel one another. For this reason Dirichlet, returning to the original concept of convergence of infinite series, investigated the limiting value reached by the sum of a number of terms when this number is assumed to increase without bound, and he completely fathomed this question by the exact determination of the limiting value of a simple definite integral, which, because of its many applications, has since been counted among the foundations of the theory of definite integrals.

By the same method and with the same means Dirichlet also carried out the more general and more complicated investigation of the convergence of the expansion of an arbitrary function of two independent variables ordered according to spherical functions. It was moreover necessary only to choose the expression of the spherical functions by definite integrals in such a suitable way that the sum of an indeterminate number of the first terms of this expansion, whose coefficients are given as double integrals, became as simple as possible and appeared in a form in which its limiting value could easily be determined by means of the limiting value already found for that simple integral.

Dirichlet found not only the special theory of these two kinds of series expansions, but also the general theory of infinite series, still unfounded in some essential points. It was indeed known that divergent series have no determined values; but the rules and conclusions valid for finite series were still transferred to infinite series in too naive a way, and no one had ever thought that even the most elementary rule according to which every algebraic sum is independent of the order of its parts might cease to be correct for sums consisting of infinitely many parts, until Dirichlet showed that there is a class of convergent series with positive and negative terms which can take other values, and even become divergent, when only the order of their terms is changed. The deeper insight thereby gained into the nature of infinite expressions also became decisive for the treatment of definite integrals; for Dirichlet's observation, applied to multiple integrals whose upper limits are infinite, showed that there is likewise a whole class of such integrals in which changes in the order of integration can produce quite different values.

Dirichlet treated the general theory of definite integrals with special predilection in his lectures. In them he united into a connected whole, by appropriate ordering and method and while excluding all external aids not belonging to this theory itself, results previously scattered as isolated items. In addition, he enriched this discipline by inventing a new and distinctive method of integration. Its main idea consists in introducing a discontinuous factor so that the limits within which the integrations have to remain can be overstepped in such a way that arbitrary other, though wider, and especially infinitely wide, limits can be taken instead of the given ones without changing the value of the integral. In applying this method to the attraction of ellipsoids and to the evaluation of a new multiple integral he also showed that, handled with skill, it is capable of giving the solutions of certain difficult problems by a simpler path than the other known integration methods.

While occupied with these analytical works, Dirichlet never ceased also to keep in mind the great problems of his favourite discipline, number theory, and to meditate on their solution. In his mind, always striving toward unity, he could not allow these two spheres of thought to stand side by side without inquiring into their inner relations, in which he sought, and in fact found, knowledge of many deeply hidden properties of numbers. His applications of analysis to number theory that arose from this differ essentially from all earlier attempts of the kind in that in them analysis is made serviceable to number theory in such a way that it no longer merely happens to throw off some isolated results for it, but must with necessity yield the solutions of certain general classes of problems of arithmetic still wholly inaccessible by other paths. These Dirichlet methods are epoch-making for number theory in a way similar to Descartes's applications of analysis to geometry. They would have to be recognized, just as analytical geometry is, as the creation of a new mathematical discipline if they extended uniformly not merely to certain classes but to all problems of number theory.

The first application that Dirichlet made of his new method concerns the very simple theorem that every arithmetic progression whose terms do not all have a common factor contains infinitely many primes. Because of its quite elementary character this theorem is of great importance in many arithmetical investigations, and in particular it had been used merely as an unproved result in Legendre's first attempted proof of the law of quadratic reciprocity. Euler's distinctive way of concluding, from the transformation of a product containing only primes into a divergent infinite series, that the number of all primes is infinite, gave Dirichlet the occasion for the more general application of infinite series and infinite products. In seeking to arrange these analytical aids in accordance with the purpose of his investigation, he arrived at the fundamental theorem of his new method, which determines the limiting value of a general series of powers of positive decreasing quantities whose common exponent approaches the limit one. The application to the proof of the theorem on arithmetic progressions requires the development of a certain group of infinite products into infinite series, and then the task is to prove that the product of these infinite series becomes infinitely large when the common power exponent of all terms approaches the limiting value one. Since the first factor of this product, according to the fundamental theorem just mentioned, necessarily becomes infinitely large, it remains only to show that the product of all the other factors cannot be zero. Although these infinite series can be summed in finite form by means of logarithms and circular arcs, the complete settlement of this important point, especially for the case where the common difference of the arithmetic progression is a composite number, presented very considerable difficulties, which Dirichlet in the first treatment of the investigation could overcome only by very complicated and indirect considerations. But precisely this difficulty became the occasion for a second application of analysis to number theory, which contains one of his most important and brilliant discoveries: the determination of the number of classes of quadratic forms for every given determinant. The difficulty of the first investigation was disposed of by this second one in the simplest way, for it showed of itself that the product in question cannot be zero unless the number of classes of quadratic forms would also have to be zero.

The determination of this class number rests on considering the doubly infinite sum whose general term is unity divided by a power of a quadratic form, the exponent being taken infinitely near to one. This sum, which extends over all integral values of the two indeterminate variables, with certain restrictions, and over all non-equivalent forms of the same determinant, is determined in two different ways: once by taking together the representative forms of all classes, and the other time by considering them individually. Since the sum receives the same value for each of these classes, the total sum is obtained as such a single sum multiplied by the number of classes. The form in which the expression for the class number thus obtained finally presents itself is essentially different for negative and for positive determinants, and in both cases reveals a surprising connection of the class number with quite different regions of number theory: for negative determinants, with quadratic residues and nonresidues; and for positive determinants, with the two solutions of Pell's equation distinguished above all others, namely the fundamental solution containing the smallest numbers and the solution arising from the theory of division of the circle, which Dirichlet first gave in a short paper on the solution of Pell's equation by circular functions. Since then no deeper insight has been gained into the connection of these apparently quite heterogeneous objects with the class number and with one another, because in general no method other than Dirichlet's exists that would be capable of solving such difficult questions.

Although the fame of this great discovery belongs to Dirichlet alone, insofar as he drew it entirely from his own mind, a certain share that Jacobi and Gauss have in it must nevertheless not be left unmentioned. From comparing certain theorems of cyclotomy and of the composition of quadratic forms, Jacobi had already, some years earlier, more guessed than derived the class number for forms whose determinant is a negative prime; and since a series of numerical examples confirmed his conjecture, he had also published it. But this could have had no influence on Dirichlet's discovery, because according to his method he could not aim to prove a determinate result in a finished form, but only to find it, and in such a way that nothing whatsoever about its form could be predetermined. Gauss, however, as the papers he left behind have shown, had long been in possession of the complete expression for the class number for negative determinants, which he had found not by induction but probably by a method similar to Dirichlet's. That he never published this result, nor a whole series of the most important and brilliant discoveries later made by Abel and Jacobi and found in his writing desk, seems to have its reason not only in an inexplicable peculiarity of this extraordinary man, but probably also in the fact that the methods he used may not have satisfied him in all points, and that he preferred to give nothing rather than something defective.

Dirichlet extended the application of his method to the theory of quadratic forms not to the class number alone but also to all the related questions concerning the division of the classes into genera and orders. In this way he made accessible by a new path the most interesting, but because of the difficulty of the methods the hardest to understand, section of Gauss's Disquisitiones Arithmeticae. Moreover, according to principles similar to those used for arithmetic progressions, he also proved for quadratic forms the theorem that every form whose three coefficients have no common factor represents infinitely many prime numbers.

When, later, in a special memoir, he treated the theory of quadratic forms from the viewpoint that the coefficients and the variables of the form are to be considered as complex integers arising from the decomposition of the sum of two squares into imaginary factors, he obtained from the application of his method to them several new and surprising results. I must especially emphasize two of them, which are of special importance because they open deeper views into the most hidden properties of forms of higher degree. The class number of quadratic forms with complex coefficients is expressed by series which, as Dirichlet unfortunately did not develop fully but only indicated, can be summed not by circular functions but by lemniscatic functions. These therefore stand here in the same relation to the solutions of the corresponding Pell equation, or more generally to the units, as the circular functions stand to the units of non-complex quadratic forms. Since the forms considered here can also be conceived as decomposable forms of the fourth degree with four variables, this result points in general to an intimate, still unexplored connection in which certain forms of higher degree stand to definite, and indeed periodic, transcendental functions of analysis. Dirichlet further found that in the special case where the determinant is a non-complex number, the class number of these complex quadratic forms consists of two factors, both of which separately express the class numbers of non-complex forms of the same determinant, one for the negative determinant and the other for the positive. This example first revealed the more general nature of these expressions, which recurs in all class numbers of forms of higher degree subsequently determined: namely that they consist of two essentially different integral factors, one determined solely by the units, the other by power residues with respect to the determinant.

For those decomposable forms of higher degree whose linear factors contain no irrationalities other than roots of unity with prime exponent, Dirichlet determined the class number during his stay in Italy, but unfortunately he published nothing of this work.

Finally, one must also mention here the interesting and new results which Dirichlet obtained from the application of his method to the determination of mean values, or asymptotic laws, for the integral functions appearing everywhere in number theory and seeming to progress quite irregularly. They concern the question, already treated earlier by Euler, Legendre, and Gauss, of the frequency with which primes occur in the natural sequence of numbers; also the mean values indicated by Gauss for the number of classes of quadratic forms and the number of their genera; and in addition several functions occurring in the elements of number theory and relating to divisors and residues of numbers. Remarkably, in precisely this type of investigation, for which analytical treatment appears especially suited, Dirichlet's continued efforts succeeded in many cases in replacing analytical methods by purely arithmetical ones, and in this way in obtaining some further new and surprising results. I shall mention only one: when a given number is divided by all smaller numbers, the remainders that are less than half the divisor occur on average much more frequently than those that are greater.

The investigations mentioned above concerning certain forms of higher degree with several variables led Dirichlet to the general theory of decomposable forms of all degrees, which is essentially identical with the general theory of the complex numbers introduced into science by Gauss. What he published on this important subject is indeed limited to a few sketch-like brief communications in the monthly reports of our Academy; nevertheless, besides the general guiding ideas, they contain the solution of the principal fundamental difficulties. In particular, the theorems on the units and the chief points in the proofs of them are stated so completely that this part of the theory could be worked out wholly in the spirit of its author. The significant benefit that would accrue to science from a complete treatment of this general theory can already be clearly recognized from what Dirichlet gave of it; for in it one finds again the most essential and beautiful properties of the relevant special theories, notably also of quadratic forms, which in the generality corresponding to their nature do not appear in more complicated shape, but, purified of the inessential determinations that adhere to everything special, can first be recognized in their true simplicity and greatness.

The lectures on number theory which Dirichlet introduced at the university here, and indeed first introduced at German universities in general, also caused him to devote special diligence to the more elementary parts of this discipline, and especially to the simplification of Gaussian methods and proofs. Among the works belonging here, which he not merely communicated orally to his students but also published elsewhere when occasion arose, I first mention the new treatments of two Gaussian proofs of the law of quadratic reciprocity: namely, the first proof given in the Disquisitiones, which even in Dirichlet's clear and appropriate treatment is inferior in simplicity to other proofs of this theorem and has only this in its favour, that it requires no aids other than those lying within the theory of quadratic residues itself; and the fourth Gaussian proof, derived from the summation of certain finite series whose absolute value can be stated very easily, while the whole difficulty lies in determining the associated sign, which Dirichlet mastered by summing these series by means of definite integrals. Further to be emphasized is the new treatment, published as an academic occasional writing, of the doctrine of composition of quadratic forms, in which he derived in the simplest way the theorems worked out in Gauss only by a hard-to-manage apparatus of formulas, by going to the essence of the matter and considering, instead of the forms, the numbers to be represented by them. The already mentioned work on the theory of quadratic forms for complex numbers can also be counted here, insofar as its simple methods are everywhere applicable to ordinary quadratic forms as well, whereby it at the same time takes the place of a simple and thorough systematic presentation of this elementary theory. Finally there belong here the new proofs of the theorems on the number of decompositions of numbers into four and into three squares, and the general reduction of positive quadratic forms with three variables, the latter first carried out by Seeber in an extremely complicated manner. In general, one recognizes in the methods by which Dirichlet simplified number theory in these works and made it more accessible that they are drawn chiefly from the thorough study of the more general theories. The proofs of the theorems therefore rest not on special and accidental determinations, but throughout on the essential properties of the number-theoretic concepts concerned; and in the special case they thereby mediate at the same time the recognition of the general.

Dirichlet's works in mathematical physics and mechanics originally proceeded from Fourier's theory of heat, which, as already remarked, was also the source of his first analytical works. He published, however, only one work concerning the theory of heat itself: a rigorous and simple solution of the problem, already treated by Fourier, of determining the heat state of an infinitely thin rod for which the temperatures at both ends are given as functions of time.

Later, the questions and mathematical-physical investigations stimulated and carried out by Gauss came to predominate in his interest; and he chose especially the theory of forces acting inversely as the squares of the distances as the object of his researches, on which he also gave special lectures at the university. Of the two memoirs he published in this field, one gives the solution of the problem of finding the density with which an infinitely thin layer of mass must be laid upon a spherical surface so that the potential at every point of it has a given continuously varying value. Gauss had shown that this problem has a determinate solution for every surface and that for the spherical surface this solution is also analytically executable. It was especially necessary here to investigate the convergence of the expression for the density developed according to spherical functions, which follows easily from the corresponding expression of the given potential value. That convergence did not follow directly from Dirichlet's above-mentioned memoir on the convergence of series ordered according to spherical functions, since the density could also be infinite at places. The exact investigation of this point gives the result that the convergence of this series in fact does not take place generally without exception, but rather is subject to certain conditions whose absence indeed causes the series to diverge. The final result is then simplified by carrying out the summations so that it requires no infinite operation other than a double integration. The second short memoir belonging here concerns the potential as such and contains in this respect a new definition of it, since Dirichlet shows that the known equation among the second partial differential quotients, together with certain conditions of continuity and finiteness satisfied by the potential and its first differential quotients, determines it in such a way that no analytic function other than the potential satisfies all these conditions. Accordingly every expression found for a potential can be tested and verified a posteriori by differentiation. This new way of defining analytic functions by means of continuity conditions has since been raised by Dirichlet's successor in Göttingen, Professor Riemann, to a special principle of analysis, which has already proved extraordinarily fruitful in his works and seems destined to found a new epoch in the direction recently pursued by analysis, in which the solution of its problems is forced less by calculation than by thought.

The investigation of the stability of equilibrium, in which Dirichlet first rigorously proved that to every maximum of the force function there really corresponds a position of stable equilibrium, was carried out in a similar spirit. By using not the analytical rules for the determination of maxima of functions but only the original concept of a maximum, it attained not only the exceptionless general validity lacking in all earlier proofs but also a wonderful simplicity and clarity.

In his investigations on the motion of fluids Dirichlet gave the first example of a truly executed integration of the general differential equations of hydrodynamics: namely, for the case in which in an infinite mass of fluid, originally at rest, a sphere moves under arbitrary accelerating forces in a constant direction and by its motion sets the fluid itself in motion. He found in this connection the very remarkable result, contradicting ordinary ideas about the resistance of fluids, that the resistance which the sphere has to suffer in its motion depends not on its velocity itself but only on the increase of that velocity; so that when the accelerating forces cease to act on the sphere, the resistance also disappears, and the sphere makes a constant motion in a straight line in the fluid.

Finally, there must be mentioned the memoir on a problem of hydrodynamics, which was also Dirichlet's last work. It gives another example, not merely approximate but rigorous, of the integration of the general hydrodynamic equations, in the determination of the motion of a fluid under the assumption that the individual particles of mass, in their motion, continually preserve a certain condition of affinity and that the original form of the fluid is that of an ellipsoid. Dirichlet proves that such a motion is indeed possible and that throughout its whole duration the fluid retains the shape of an ellipsoid, with the same centre but with variable position and magnitude of the principal axes. In the special case where the ellipsoid is one of revolution, all integrations can be completely reduced to quadratures, and the fluid makes isochronous oscillations, alternately taking the form of a prolate and an oblate ellipsoid.

Before I now return from considering Dirichlet's scientific works to describing his life, I must emphasize a general observation suggested by comparing them with Jacobi's works. Since these two men worked contemporaneously for a quarter-century at the further development of the mathematical sciences, and stood in close personal friendship and active scientific intercourse with one another, it is a very striking phenomenon that their writings, although they often concern the same special disciplines, show almost no immediate points of contact. The special objects of their researches are, with very few insignificant exceptions, entirely different, and there are scarcely any examples even of one using the results of the other for his own investigations. This lack of relation in their writings cannot be sufficiently explained solely by the difference of the starting-points and directions of their mathematical studies and labours; rather it is grounded in the fact that both deliberately avoided encroaching upon those regions in which each recognized the superiority of the other, and that they sought to avoid even the appearance of rivalry.

The first personal acquaintance between Dirichlet and Jacobi was formed in 1829, when Jacobi came from Königsberg to Berlin to visit his relatives and friends. On a journey they undertook together to Halle, and from there in the company of Mr. W. Weber to Thuringia, they became better acquainted. Since Jacobi often spent his vacation time in Berlin, they later had opportunity for more intimate scientific and friendly intercourse. Afterwards, when Jacobi was seized by a dangerous illness and, on the advice of physicians, had to seek the milder climate of Italy for his recovery, Dirichlet, who had long intended a journey to Italy, seized this opportunity to spend a winter in Rome together with Jacobi, and in the autumn of 1843 set out there with his whole family. Since at the same time our colleagues Mr. Steiner and Mr. Borchardt also spent that winter in Rome, German mathematics was at that time represented there in a very brilliant and many-sided fashion. Dirichlet remained in Italy a year and a half, extended his journey also to Sicily, and spent the next winter in Florence. On his return he found Jacobi in Berlin, for in the meantime, by the grace and munificence of His Majesty the King, Jacobi had been given leave from Königsberg and called here, so that, without occupying a definite office, he could care for his health and live wholly for science.

The common interest in knowing truth and promoting the mathematical sciences remained the firm foundation of the friendly relation in which Jacobi and Dirichlet lived together here. They saw one another almost daily and discussed with one another more general or more special scientific questions, whose spirited treatment, precisely because of the difference of the standpoints from which both surveyed the whole field of the mathematical sciences, always retained a new and living interest. Jacobi, who through the wonderful abundance of his mind no less than through the depth of his mathematical researches and the brilliance of his discoveries knew how to secure everywhere the recognition due him, enjoyed at that time a much more widespread reputation than Dirichlet, who did not possess the art of making himself count, and whose writings, chiefly treating only the most difficult problems of science, had a less extensive circle of readers and admirers. No one recognized this disproportion between Dirichlet's outward recognition and his scientific importance more correctly than Jacobi; and no one else was at the same time more skilful and more active in equalizing it and in procuring for his friend, even in wider circles, the recognition he deserved. It is chiefly to Jacobi's activity also that Dirichlet was preserved for our Academy when, in 1846, the government of Baden intended to win him for the University of Heidelberg. Two letters which he addressed in this matter to Alexander von Humboldt and to His Majesty the King give, in a few strong and apt strokes, a vivid presentation of Dirichlet's scientific greatness and of the irreplaceable loss that the exact sciences in Prussia, the Academy, the University, and especially Jacobi himself, would suffer if Dirichlet were to leave our fatherland.

This threatening loss, happily averted at that time, confronted us nine years later, after Jacobi and Gauss had passed away, all the more painfully. The University of Göttingen, which for half a century had enjoyed the fame of possessing the first among all living mathematicians, strove eagerly to maintain this fame further by calling Dirichlet to Gauss's position, and first addressed to him the question whether, and under what conditions, he might be inclined to accept such a call. Dirichlet had here a sphere of activity such as he could hardly expect to find again at another university; he enjoyed in high degree the veneration of his listeners and the respect and affection of his colleagues, and besides was bound to Berlin by close family ties. The only thing that made a change in his situation desirable to him was that teaching at the War School fragmented the powers that he would gladly have devoted wholly to the university and to science. It was therefore his lively wish to be released from the position at the War School and to have this loss of income covered by the university. Since the call to Göttingen offered him the opportunity to attain his purpose in one way or another, he declared in answer to the inquiry addressed to him that he would obey an official call from the Royal Hanoverian government unless, before its arrival, his position here were changed in accordance with his wishes. His friends, to whom he communicated this, did not omit to inform the Royal Ministry of it, so that timely precautions might be taken to avert the threatened loss from the university here and from the Academy. But Minister von Raumer did not wish to make an immediate decision; he wanted first to await an official step by the Royal Hanoverian government. That government soon sent Dirichlet his formal call through his friend Professor Weber in Göttingen, who delivered it personally. When now Minister von Raumer, in order to keep him here, offered him even more than he had wished, it was too late. Since Dirichlet now considered himself bound by his previous declaration, no advantages or considerations could determine him otherwise.

In the autumn of 1855 he moved from here to Göttingen. There he arranged for himself, entirely to his liking, his own pleasantly situated house with a garden. The quiet of the smaller town, which he had not enjoyed since his youth, compensated him sufficiently for the outward amenities of metropolitan life in Berlin. He also found there men of like mind to whom he could attach himself more closely; and his scientific and general intellectual importance, combined with the unpretentiousness and honourableness of his whole nature, soon won him the same general esteem that he had enjoyed here. At the university he did not indeed find as large a circle of listeners as the one he had left here; but his reputation as a teacher, acknowledged no less than his scientific reputation, drew to Göttingen many young men striving for higher training in the mathematical sciences, and some of the most distinguished academic teachers there also became his eager listeners, so that the success of his lectures there was in proportion no less than here. Since his mathematical researches, which always lay closest to his heart, were also favoured by the greater leisure he enjoyed, he felt very satisfied in his new position.

In the summer of 1858, after the close of his lectures, he travelled to Switzerland and stayed in Montreux on Lake Geneva, less for rest than to work out there a memorial address on Gauss to be delivered in the Göttingen Society of Sciences, and a memoir for that society's transactions. When he had almost completed the hydrodynamical memoir mentioned above, he was suddenly seized by an acute heart disease and hastened at once back to his family in Göttingen, where he arrived mortally ill. The art of the physicians and the loving care of his family succeeded in happily averting the immediate danger to his life; but he could scarcely raise himself again somewhat from his sickbed, and for his hoped-for complete recovery still needed the greatest rest of body and mind, when his wife, suddenly struck by apoplexy, died after a few hours, without his being able to see her once more. This unexpected blow turned his illness again for the worse, and after severe sufferings he succumbed to it on 5 May 1859.

As a human being Dirichlet was no less distinguished by his noble character than in science by the depth and solidity of his mind. The honourableness that filled his whole nature and appeared pure and unclouded in all his actions sprang from the high moral cultivation of his mind and heart, and for that reason was directed not to outward honour but everywhere only to true inner honour, whose exact and strict measure he bore within himself. Ambition, which seeks outward recognition and finds its satisfaction more in appearance than in essence, was wholly foreign to him. Even the scientific honours from learned bodies that were bestowed upon him in the richest and highest measure he valued chiefly insofar as he could see in them the approval of connoisseurs and competent judges; they had no influence on the clear judgment which he exercised with complete impartiality concerning the value of his own achievements.

As in science, so also in his whole life the love of truth was the moral foundation of his thinking and acting. It pressed back the activity of imagination in him, kept him free from prejudices and self-deceptions, and allowed him full satisfaction only where he could arrive at exact and completely secure knowledge. Truth in symbolic form corresponded less to his nature; truths that announced themselves as results of philosophical speculation appeared to him in general suspect. He used to say of philosophy that it was an essential defect of it that it had no unsolved problems, as mathematics does; that it therefore was conscious of no definite boundary within which it had really investigated truth and beyond which, for the time being, it must modestly confess that it knew nothing. The greater the claims philosophy made to omniscience, the less perfectly clear recognized truth he believed he could entrust to it, since from his own experience in the field of his science he knew how difficult the knowledge of truth is and what effort and labour it costs to advance it even one step.

A certain shyness that had belonged to Dirichlet in his youth had, in maturer age, ennobled itself into true inner modesty; yet it still showed itself in many respects as natural reserve, especially in the fact that he very unwillingly appeared in public, did not like to take the floor in larger assemblies, and never delivered speeches unless it was for him an unavoidable duty. In general he never thrust himself forward, neither with his person nor with his opinions and judgments, but was reserved even where his judgment as an expert was requested, because precisely in such cases he proceeded with the greatest conscientiousness and gave his definite judgment only after considering all sides. To his disposition, directed more toward knowledge than toward practical activity, every striving after outward influence was foreign. In his outward relations in life he in fact never made any influence effective except that which a noble and spirited man involuntarily and immediately exercises in the circles to which he belongs.

In social and friendly intercourse Dirichlet everywhere showed the genuine humanity that is grounded in general respect for the personality of human beings and in freely allowing their peculiarities and convictions. He had an open eye for the good sides of others and preferred to seek these out rather than to dwell on their weaknesses and defects, which he never made the object of self-satisfied mockery and opposed only when they betrayed a lack of honourable disposition. The same humanity he showed in his extensive scientific intercourse with the most important and capable mathematicians at home and abroad, an intercourse he preferred to maintain personally rather than by letter, because letter-writing was not pleasant to him, while on journeys he gladly visited his acquaintances and was often visited by them. He always showed a very lively participation in the achievements of others, gladly entered in conversation into their special scientific interests, and instructed by communicating the higher points of view from which he surveyed the questions at hand, without ever making the superiority of his mind be felt in an oppressive way.

The most capable among the younger German mathematicians had almost all formerly been Dirichlet's listeners, and esteemed him not merely as their teacher, to whom they owed the best part of their mathematical education, but were also always attached to him with true love and veneration. How highly he for his part knew how to value the love of his pupils, and how above all he recognized it as the highest reward of his teaching activity, he expressed shortly before his death in a beautiful and worthy way. After one of the last severe attacks of his illness, when he felt somewhat freer again, he expressed the wish to see once more one of his dearest friends and former pupils. The latter, being informed of this, immediately travelled to him and had the good fortune, during two days in which the illness had somewhat abated, to see and speak with his beloved teacher. At parting Dirichlet said to him: It is truly rewarding to be a professor when one earns such love.

The success of his teaching activity, externally measured by the number of listeners, especially in the later period of his academic work, was as considerable as probably no teacher at a German university can show in the field of higher mathematics. He owed it to no didactic artifices, nor to the gift of a brilliant delivery, but solely to the inner clarity of his mind, by virtue of which he knew how to grasp and present even the most difficult subjects in their simple truth. In doing so he spared his listeners no exertion of thought necessary for complete knowledge of the subject, but he gladly spared them and himself lengthy and time-consuming calculations by replacing them, wherever possible, by simple thoughts. If one measures the success of his teaching by the scientific competence of the younger mathematicians who were his students and who owe chiefly to him their mathematical education, then only Jacobi's activity can in general be regarded as equal to his, and perhaps in one respect even more highly valued: Jacobi founded a special mathematical school that continues to work in his spirit and sense, whereas Dirichlet's students pursue more individually different directions.

His own scientific direction was so closely connected with the peculiarity of his mind and character that it could not become the common property of a school. He did not love the much-trodden and already levelled roads of science; rather he took pleasure in choosing as the object of his reflection and work the principled difficulties that those roads usually avoid. When he had fathomed them, he did not linger in spinning out the consequences of the results obtained, but preferred to work on from them into the depths, where he found new difficulties to overcome. For this reason his writings are not extensive and consist mostly of smaller memoirs in which he treats definite problems of science and completely fathoms them. Especially characteristic of his scientific direction is also the perfect rigour and evidence of the methods and proofs by which he establishes his results. This property, although it corresponds only to a demand lying in the essence of mathematics itself, is nevertheless found even among the greatest mathematicians only rarely in complete purity. In the field of analysis it first came into force through Gauss, and since then has so little become common property that even Jacobi's writings show its absence at certain places, as Jacobi himself openly admitted.

That Dirichlet kept himself and his writings free from such defects he owed chiefly to the love of pure and perfectly secure truth that was proper to him, and in addition to the way in which he worked and to the care with which he composed his writings. The clarity and definiteness of his thinking and the unusual power of his memory, by virtue of which he kept once thought and investigated matters perfectly present to himself at any time, made the use of the pen while working almost entirely dispensable to him. Nor did he need special quiet or leisure for it; on walks, on journeys, at musical entertainments, and generally in all situations where he did not himself need to speak or act, he could continue his deep speculations with the same success as at his writing desk. As an example I may mention that he fathomed the solution of a difficult problem of number theory, with which he had struggled in vain for a long time, in the Sistine Chapel in Rome while listening to the Easter music usually performed there. When he had found important results, he devoted the greatest diligence to finding, through investigation of their connections on all sides with one another and with related theorems, the simplest method of derivation most appropriate to the nature of the subject. Only after he had succeeded in this did he proceed to written elaboration, to which he usually resolved with difficulty, but which he then carried out with the greatest care.

Of the results that Dirichlet worked out in the last years of his life, only little has been preserved for science, because he postponed their written elaboration too long. Among the papers he left behind there was found no mathematical manuscript except the one hydrodynamical memoir that recently appeared in the Transactions of the Göttingen Society, edited by Professor Dedekind, to whom Dirichlet himself had still entrusted its completion. From what he occasionally communicated to individual friends about the subjects of his research, it appears that among other things he had completed in his head a full theory of ternary indeterminate forms of the second degree; that he had succeeded in carrying the approximation of the asymptotic laws for a kind of number-theoretic functions, on which the determination of the frequency of primes depends, a whole degree further; and that he had found a mathematically completely rigorous proof of the stability of the world system. Of a great and especially valuable discovery from the last period of his life, namely a wholly new general method for the treatment and solution of problems of mechanics, he spoke only once in the summer of 1858 to one of his friends, Mr. Kronecker, with whom he stood in the most intimate scientific and friendly intercourse. He himself attached quite special weight to this discovery and asked Mr. Kronecker for the time being to speak of it to no one. Kronecker therefore communicated to his friends only after Dirichlet's death what he had learned from him about it: namely, that this method was not directed toward reducing the integrations of the relevant differential equations to quadratures, because this means, by which Jacobi had attempted to obtain the solution of mechanical problems, was too limited; that his procedure consisted rather in a stepwise approximation, in which each new step at the same time afforded a more complete and more exact insight into the nature of the motions determined by the conditions of the problem; and finally that the theory of small oscillations offered a certain support for finding this method.

I will give words to the lament over these losses to science, perhaps not to be replaced for a long time and whose magnitude can be sufficiently measured from the indications still available, only by recalling the saying which Dirichlet himself uttered in the memorial address on Jacobi concerning that man's unfinished works: Death, which took him away from his labour too early, did not grant such great enrichments to science.
\end{document}
